\VerdierPhysicalPageAnchor{0031}%
\VerdierSemanticAnchor{intro-8-etale-lefschetz}
In the setting of the étale topology, there is an analogous Lefschetz formula
[12].
\VerdierSemanticAnchor{intro-8-frobenius-L}
Applied to the Frobenius endomorphism of a smooth proper algebraic curve over
a finite field, equipped with a sheaf constructible for the étale topology,
this formula made possible the proof of the rationality of generalized
$L$-functions [11]. Note that this proof uses a Lefschetz formula in which
the coefficient sheaf does indeed have singularities.

\VerdierSemanticAnchor{intro-9}
\noindent\textbf{9.}\quad
In the case of the topological Lefschetz formula, we saw that the complex
\VerdierSemanticAnchor{intro-9-perfect-complex}
\[
\mathrm{RH}(X,F)
\]
is perfect, which makes it possible to define the Lefschetz number of
\VerdierSemanticAnchor{intro-9-rh-phi}
\[
\mathrm{RH}(\Phi)\colon
\mathrm{RH}(X,F)\longrightarrow\mathrm{RH}(X,F).
\]

\VerdierSemanticAnchor{intro-9-cohomology-objects}
This shows in a particularly striking way why a complex cannot be replaced
by the collection of its cohomology objects. Indeed, the trace of the
endomorphisms
\VerdierSemanticAnchor{intro-9-hi-phi}
\[
\mathrm{H}^i(\Phi)\colon
\mathrm{H}^i(X,F)\longrightarrow\mathrm{H}^i(X,F)
\]
can be defined only under extremely restrictive conditions.

\VerdierSemanticAnchor{intro-9-regular-trace}
When, for example, the ring $A$ has finite global cohomological dimension,
i.e. when $A$ is regular, one can define the trace of $\mathrm{H}^i(\Phi)$
by taking finite projective resolutions. One then has
\VerdierSemanticAnchor{intro-9-trace-formula}
\[
\chi\bigl(\mathrm{RH}(\Phi)\bigr)
=\sum_i(-1)^i\operatorname{Tr}\bigl(\mathrm{H}^i(\Phi)\bigr).
\]

\VerdierSemanticAnchor{intro-9-nonregular-rings}
But, as we have seen, particularly in sheaf theory for the étale topology,
one has to work with rings of the form
$\mathbb{Z}/p^\ell\mathbb{Z}$ ($\ell\geq 2$), which are not regular.

\VerdierSemanticAnchor{intro-9-k-theory}
This notion of a perfect complex is also useful in algebraic geometry:
K-theory and the Riemann--Roch theorem [14].

\VerdierSemanticAnchor{intro-10}
\noindent\textbf{10.}\quad
In conclusion, it seems that at present the language of derived categories
is the indispensable tool for formulating and proving the essential results
of the theories just mentioned.

\VerdierSemanticAnchor{intro-10-additive-scope}
We have restricted ourselves in this work to additive homological algebra:
abelian categories and additive functors. As soon as one turns to the
Riemann--Roch theorems [14], this framework is inadequate.
\VerdierSemanticAnchor{intro-10-omissions}
Nor have we addressed the study of the various class groups, or Grothendieck
groups, that can be formed from triangulated categories. Finally, we have
set aside questions of relative homology, which can also be treated in the
framework of triangulated categories [15].

